\section{Определение дробно-линейной функции. Отображение прямых и
окружностей при дробно-линейном отображении. Выражение для
дробно-линейной функции, отображающей три заданные точки комплексной
плоскости в три заданные
точки.}\label{ux43eux43fux440ux435ux434ux435ux43bux435ux43dux438ux435-ux434ux440ux43eux431ux43dux43e-ux43bux438ux43dux435ux439ux43dux43eux439-ux444ux443ux43dux43aux446ux438ux438.-ux43eux442ux43eux431ux440ux430ux436ux435ux43dux438ux435-ux43fux440ux44fux43cux44bux445-ux438-ux43eux43aux440ux443ux436ux43dux43eux441ux442ux435ux439-ux43fux440ux438-ux434ux440ux43eux431ux43dux43e-ux43bux438ux43dux435ux439ux43dux43eux43c-ux43eux442ux43eux431ux440ux430ux436ux435ux43dux438ux438.-ux432ux44bux440ux430ux436ux435ux43dux438ux435-ux434ux43bux44f-ux434ux440ux43eux431ux43dux43e-ux43bux438ux43dux435ux439ux43dux43eux439-ux444ux443ux43dux43aux446ux438ux438-ux43eux442ux43eux431ux440ux430ux436ux430ux44eux449ux435ux439-ux442ux440ux438-ux437ux430ux434ux430ux43dux43dux44bux435-ux442ux43eux447ux43aux438-ux43aux43eux43cux43fux43bux435ux43aux441ux43dux43eux439-ux43fux43bux43eux441ux43aux43eux441ux442ux438-ux432-ux442ux440ux438-ux437ux430ux434ux430ux43dux43dux44bux435-ux442ux43eux447ux43aux438.}

\textbf{Определение:} \emph{Дробно-линейной} называется функция \[
w= \frac{az+b}{cz+d},\quad ad-bc\neq 0,\quad(1)
\] где \(a,b,c,d\) - заданные комплексные числа. Условие \(ad-bc\neq 0\)
означает, что \(w\not\equiv \mathrm{const}\). Отображение,
осуществляемое функцией (1) называется \emph{дробно-линейным}. При этом
предполагается, что если \(c\neq 0\), то
\(w\left( -\dfrac{d}{c} \right)=\infty,w(\infty)= \dfrac{a}{c}\), а если
\(c=0\), то \(w(\infty)=\infty\). В частности, если \(c=0\), то функция
(1) является линейной, а отображение, осуществляемое линейной функцией,
называется \emph{линейным}.

\textbf{Теорема:} При дробно-линейном отображении образом любой
окружности или прямой является окружность или прямая.

\textbf{Теорема:} Существует единственное дробно-линейное отображение,
при котором три различные точки \(z_{1},z_{2},z_{3}\) переходят
соответственно в три различные точки \(w_{1},w_{2},w_{3}\). Это
отображение определяется формулой \[
\frac{w-w_{1}}{w-w_{2}}\cdot \frac{w_{3}-w_{2}}{w_{3}-w_{1}}= \frac{z-z_{1}}{z-z_{2}}\cdot \frac{z_{3}-z_{2}}{z_{3}-z_{1}}
\]
